\documentclass[11pt]{article}
\usepackage[T1]{fontenc}
\usepackage[utf8]{inputenc}
\usepackage{lmodern}
\usepackage[margin=1in]{geometry}
\usepackage{amsmath,amssymb,amsthm,mathtools}
\usepackage{hyperref}
\usepackage{enumitem}
\usepackage{verbatim}
\newcommand{\OPHPaperReleaseID}{r1412}
\newcommand{\OPHPaperReleaseDate}{May 4, 2026}
\newcommand{\OPHPaperReleaseBanner}{%
\begin{center}
\small\textbf{Paper release:} \texttt{\OPHPaperReleaseID}\quad\textbf{Released:} \OPHPaperReleaseDate
\end{center}
}


\hypersetup{colorlinks=true,linkcolor=blue,citecolor=blue,urlcolor=blue}

\title{Screen Microphysics, Patches, and Observer Synchronization in OPH}
\author{B.\ M\"uller}
\date{\OPHPaperReleaseDate}

\begin{document}
\maketitle
\OPHPaperReleaseBanner

\begin{abstract}
This paper is a phase-1 architecture note for Observer-Patch Holography (OPH). The
objective is constructive and simulator-facing: choose one finite local screen model that can
actually be encoded in qubits, make the map from local registers to patch algebras and overlap
observables explicit, and state how stable records, observer patterns, local interaction, and
observer-level synchronization could be built inside that same microscopic system. The working
choice is a constrained finite gauge-register or quantum-link style model on a cellulated
screen, used here as a \emph{reference architecture} rather than as a claim that the unique OPH
UV completion has already been identified. We separate reused OPH structure, architecture-level
claims, simulator targets, and conjectural bridges. In particular, this note does \emph{not}
claim that the OPH continuum or gravity branch has already been derived from the present
microphysics. The deliverable of this round is a sharper blueprint: local Hilbert spaces,
qudit-to-qubit encoding agenda, patch and overlap observable dictionary, observer criterion,
synchronization and repair options, a first implementable reference model, and a validation
package that can fail.
\end{abstract}

\tableofcontents
\newpage

\section{Project Charter}

\textbf{Status.} Extra side paper. Not part of the compact recovered core. Not yet part of the
release-critical OPH theorem stack.

\textbf{Purpose.}
\begin{enumerate}[leftmargin=2em]
\item Give a concrete reference architecture for local finite screen degrees of freedom.
\item Make the bridge from local registers to patches, overlaps, records, observers, and
synchronization explicit enough to build toy simulators.
\item Keep theorem claims, simulator claims, and conjectural bridges sharply separated.
\item Feed the existing OPH microphysics, repair, and observer task lanes without pretending
to close them automatically.
\end{enumerate}

\textbf{Non-goals.}
\begin{enumerate}[leftmargin=2em]
\item Do not claim that the unique UV completion is known.
\item Do not claim that the OPH scaling-limit continuum or gravity branch is already derived
here.
\item Do not treat ``observer'' as a consciousness claim.
\item Do not treat synchronization of overlap data as a proof of full continuum closure.
\end{enumerate}

\textbf{Phase-1 output.} A reference microscopic screen model, a qubit-encoding agenda, a patch
and overlap dictionary, an observer criterion, a synchronization design space, and a simulator
validation package.

\section{Claim Taxonomy}

The point of this paper is not to maximize the number of statements called ``derived.'' The
point is to keep the layers clean.

\begin{enumerate}[leftmargin=2em,label=\textbf{T\arabic*.}]
\item \textbf{Reused OPH structure.} The OPH screen-net picture, overlap consistency language,
local MaxEnt and refinement-stability branch, recoverable generalized entropy package, and the
regulator-scale edge-center completion language are reused as background structure. This paper
does not re-prove them.

\item \textbf{Architecture claims.} Once one finite gauge-register reference architecture is
chosen, the patch algebra, overlap algebra, edge-sector observables, record layer, and
observer-level interfaces can be defined explicitly at fixed cutoff.

\item \textbf{Simulator claims.} Small-lattice digital or tensor-network simulations can test
whether the reference architecture exhibits gauge-invariant patch observables, stable records,
low conditional mutual information across collars, and robust synchronization under local repair
schedules.

\item \textbf{Conjectural bridge claims.} It is plausible that an appropriate refinement-stable
phase of such models lands in the OPH modular-geometric branch and later supports the OPH
continuum or gravity story, but that bridge is not proved here.

\item \textbf{Excluded claims.} This note does not identify the final OPH microphysics, does not
establish a unique repair law, and does not close the Standard Model or continuum branches.
\end{enumerate}

\section{Why a Finite Gauge-Register Reference Architecture}

Several microscopic families are logically possible: constrained finite gauge registers,
quantum-link truncations of compact gauge systems, string-net or quantum-double models,
generic tensor-network ans\"atze, or ad hoc observer automata attached to a screen. For a
phase-1 constructive draft, the first working model should be selected by engineering
criteria rather than by aesthetic breadth.

The selection criteria are:
\begin{enumerate}[leftmargin=2em]
\item finite local Hilbert spaces at fixed cutoff;
\item exact local constraints rather than only emergent ones;
\item a natural patch and overlap algebra;
\item explicit edge-sector or center observables on collars;
\item straightforward qudit-to-qubit compilation;
\item a path from local mismatch syndromes to repair or synchronization maps;
\item extensibility from toy groups such as $\mathbb Z_2$ or $S_3$ to richer quantum-link
truncations.
\end{enumerate}

A constrained finite gauge-register or quantum-link style screen model is therefore the
default reference architecture for this draft. It is finite-dimensional, local, gauge-aware,
and already speaks the same language as the OPH regulator package. It also makes the collar
center and edge-sector story concrete instead of metaphorical.

This choice is a working \emph{reference architecture}. It is not a claim that all other
microscopic families are ruled out.

\section{Reference Architecture: Local Screen Registers on a Cellulated Sphere}

Take a finite cellulation $\Gamma=(V,E,F)$ of a horizon screen $S^2$. Each oriented link
$e\in E$ carries a finite register $\mathcal H_e$. The simplest version is a finite-group
register with basis $\{|g\rangle_e:g\in G\}$ for a finite group $G$, so
\begin{equation}
\mathcal H_e \cong \mathbb C^{d_G},\qquad d_G=|G|.
\end{equation}
A more flexible version replaces $|g\rangle$ by a truncated quantum-link basis labeled by a
small irrep set and matrix indices. For phase 1, the finite-group presentation is easier to
simulate, while the quantum-link presentation shows how the same pattern could extend to
compact-group truncations.

To support records and observer logic, attach a small record register to selected vertices or
coarse cells:
\begin{equation}
\mathcal H_v^{\mathrm{rec}} \cong (\mathbb C^2)^{\otimes m_v}.
\end{equation}
The extended Hilbert space is then
\begin{equation}
\widetilde{\mathcal H}_{\Gamma}=
\bigotimes_{e\in E}\mathcal H_e
\otimes
\bigotimes_{v\in V}\mathcal H_v^{\mathrm{rec}}.
\end{equation}

At each vertex $v$, a local gauge action acts on the incident oriented links by left/right
multiplication or by the corresponding truncated quantum-link action. Let $P_v$ denote the
Gauss projector onto the locally admissible subspace. The physical cutoff Hilbert space is
\begin{equation}
\mathcal H_{\mathrm{phys}} = \Bigl(\prod_{v\in V}P_v\Bigr)\widetilde{\mathcal H}_{\Gamma}.
\end{equation}

A minimal reference Hamiltonian or Trotterized circuit family should contain four layers:
\begin{equation}
H_{\mathrm{ref}} =
\lambda_G\sum_{v}(1-P_v)
+\lambda_B\sum_{f}(1-B_f)
+\lambda_R\sum_{a}(1-R_a)
+H_{\mathrm{sync}},
\end{equation}
where $B_f$ are plaquette or face operators, $R_a$ are record-stabilization or record-write
terms, and $H_{\mathrm{sync}}$ is a local overlap-repair interaction to be chosen from a small
menu later in the paper. The exact coefficients and update schedules are not fixed here.
What matters is the architecture: gauge constraints, face dynamics, record dynamics, and
repair dynamics all live on one finite screen.

\begin{verbatim}
Microscopic screen picture

   record q(v1)      record q(v2)      record q(v3)
        o------------------o------------------o
        | \              / | \              / |
        |  \    f1      /  |  \    f2      /  |
        |   \          /   |   \          /   |
        |    \        /    |    \        /    |
        o-----\------/-----o-----\------/-----o
      q(v4)    \    /    q(v5)    \    /    q(v6)
                 \  /                \/
                  \/                 /\
                  /\                /  \
                 /  \              / f4 \
                / f3 \            /      \
               /      \          /        \
              o--------\--------o----------o
          record q(v7)          q(v8)     q(v9)

   links carry finite gauge registers
   vertices host Gauss checks and optional record qubits
   faces carry plaquette/holonomy dynamics
\end{verbatim}

The simulator-facing question is therefore not ``what is the final UV completion?'' It is
``can one finite local screen model realize the OPH patch, edge, record, and synchronization
interfaces well enough to be testable?''

\section{Qudit-to-Qubit Encoding Agenda}

If the screen model is ever going to be built or simulated on conventional hardware, the
local finite registers must be compiled into qubits. This paper does not pick a single global
encoding forever; it sets the agenda and the tradeoffs.

\subsection{Binary, unary, and factorized encodings}

For a $d$-level link register, the default binary encoding uses
\begin{equation}
 n_e = \lceil \log_2 d\rceil
\end{equation}
qubits per link, with a penalty term suppressing unused binary strings when $d$ is not a
power of two. Binary encoding is qubit-efficient and natural for digital simulation.

Unary or one-hot encoding uses $d$ qubits per link with one excitation allowed. It is less
qubit-efficient but makes local projectors and permutation actions shallower. This is useful
for very small groups and for verification runs.

For truncated quantum-link bases labeled by irreps and matrix indices, a factorized encoding
is often cleaner:
\begin{equation}
|R,m,n\rangle \longmapsto |R\rangle_{\mathrm{irrep}}\otimes |m\rangle_{\mathrm{left}}
\otimes |n\rangle_{\mathrm{right}}.
\end{equation}
This keeps the representation label explicit, which is attractive when collar sector labels
are later measured.

\subsection{Agenda for phase 1}

The encoding work should proceed in the following order.
\begin{enumerate}[leftmargin=2em]
\item Start with $G=\mathbb Z_2$ or $\mathbb Z_N$ in binary encoding to validate the patch,
overlap, and synchronization pipeline cheaply.
\item Move to $S_3$ or another small nonabelian finite group to validate nonabelian sector
projectors, collar labels, and character-based observables.
\item Add a truncated quantum-link variant only after the finite-group version produces stable
measurements and sensible scaling diagnostics.
\item Compare binary and unary encodings only at the level of compilation cost and numerical
stability, not as a statement about final physics.
\end{enumerate}

What must be compiled into qubit circuits is now explicit:
\begin{enumerate}[leftmargin=2em]
\item local gauge transformations and Gauss projectors;
\item plaquette or loop operators;
\item record-write and record-read channels;
\item overlap projectors, edge-sector projectors, and mismatch syndromes;
\item local repair maps that preserve gauge admissibility.
\end{enumerate}

\section{From Local Registers to Patches and Overlaps}

The bridge from local finite registers to OPH patch language has to be explicit. A patch is
not only a geometric region. It is a connected region together with the gauge-invariant
observable algebra available there.

Let $P\subseteq \Gamma$ be a connected collection of faces, together with all vertices and
links incident to those faces. Its extended cutoff Hilbert space is
\begin{equation}
\widetilde{\mathcal H}_P=
\bigotimes_{e\subset P}\mathcal H_e
\otimes
\bigotimes_{v\subset P}\mathcal H_v^{\mathrm{rec}},
\end{equation}
and its extended local algebra is
\begin{equation}
\widetilde{\mathcal A}(P)=\mathcal B(\widetilde{\mathcal H}_P).
\end{equation}
The physical patch algebra is the boundary-fixed subalgebra
\begin{equation}
\mathcal A(P)=\widetilde{\mathcal A}(P)^{G_{\partial P}},
\end{equation}
in the same spirit as the regulator picture already used in OPH.

If $P$ and $Q$ overlap, the shared region $B=P\cap Q$ is not merely a geometric set. It is
where the shared or comparable observables live. At fixed cutoff, the collar Hilbert space is
expected to decompose as
\begin{equation}
\mathcal H_B \cong \bigoplus_{\alpha}
\bigl(H_{b_L^{\alpha}}\otimes H_{b_R^{\alpha}}\bigr),
\end{equation}
with center projectors $\Pi_{\alpha}$ labeling edge sectors. That decomposition is not proved
again here; it is reused OPH regulator language and serves as the observable interface for the
reference model.

\subsection{Patch and overlap observable dictionary}

The working dictionary is the core of this paper.

\begin{enumerate}[leftmargin=2em]
\item \textbf{Link register basis / electric occupation.}
The local computational basis on a link is the microscopic screen data. In a finite-group
presentation this is the basis $|g\rangle_e$; in a quantum-link truncation it is the local
representation basis.

\item \textbf{Gauss projector $P_v$.}
This is the local admissibility check. It says which neighboring link configurations may be
regarded as physical around a vertex.

\item \textbf{Plaquette operator $B_f$ or loop operator $W_\gamma$.}
This is the patch-level probe of curvature, flux, or sector excitation. Closed loops contained
inside a patch belong to $\mathcal A(P)$.

\item \textbf{Boundary or collar sector projectors $\Pi_{\alpha}$.}
These are the overlap-accessible center observables. They are the natural candidates for
shareable edge data and for robust cross-patch records.

\item \textbf{Record observables $Z^{\mathrm{rec}}_{v,a}$ and pointer subalgebras.}
These are commuting or approximately commuting observables stabilized long enough to act as
memories. A record is useful only if it can be read repeatedly with low back-action on the
shared sector data.

\item \textbf{Patch state $\rho_P$.}
The reduced state on $\mathcal A(P)$ is the patch description accessible to an internal
pattern supported on $P$.

\item \textbf{Overlap mismatch syndrome $S_{PQ}$.}
This is a function of shared observables, typically of the collar sector labels and the record
comparisons on $B=P\cap Q$. It is the direct input to any repair or synchronization update.
\end{enumerate}

\begin{verbatim}
Patch/overlap interface

   Patch P bulk             shared collar B            Patch Q bulk
  ----------------         =================         ----------------
  | loops W_gamma |        || Pi_alpha     ||        | loops W_gamma |
  | records R_P   | <----> || shared bits  || <----> | records R_Q   |
  | local state   |        || syndrome S   ||        | local state   |
  ----------------         =================         ----------------

   what is compared on B is not the whole quantum state
   it is a selected shared observable family on the overlap
\end{verbatim}

This is where the bridge to observers becomes concrete. An observer never needs access to the
full global state. It needs persistent access to some patch algebra and to a stable shared
observable family on overlaps with neighboring patches.

\section{Records and the Observer Criterion}

The word ``observer'' should be used operationally. In OPH language, an observer is an
internal pattern with patch access and records. In a finite screen model, that can be stated
more sharply.

\subsection{Records}

A \emph{record layer} is a commuting or approximately commuting subalgebra of robust pointer
observables, supported partly inside a patch and partly on the boundary sector interface. A
useful record variable must satisfy three practical conditions:
\begin{enumerate}[leftmargin=2em]
\item it is written by local interactions with nearby screen registers;
\item it survives for many local update cycles compared with the readout timescale;
\item it can be compared across overlaps through shared observables or recoverable summaries.
\end{enumerate}

This prevents the word ``record'' from collapsing into an arbitrary entangled correlation.
The role of the record layer is to produce stable, repeatedly accessible data.

\subsection{Working observer criterion}

The OPH manuscript already uses the tuple
\begin{equation}
O=(P_O,\mathcal A(P_O),\rho_O,R_O)
\end{equation}
for an observer patch, local algebra, local state, and records. For the reference
microphysics we add an explicit update interface and treat an observer as a pattern class
rather than as a primitive object:
\begin{equation}
O_{\mathrm{ref}}=(P_O,\mathcal A(P_O),\rho_O,R_O,\mathcal U_O).
\end{equation}
Here $\mathcal U_O$ is a local channel or update family that writes, compares, and repairs
records using only observables available in $P_O$ and on its overlaps.

A subsystem or pattern counts as an \emph{observer in the present operational sense} if it
satisfies all of the following.
\begin{enumerate}[leftmargin=2em]
\item \textbf{Patch access.} It is supported on a connected patch $P_O$ with a nontrivial
shared overlap algebra to neighboring patches.
\item \textbf{Record support.} It carries a set $R_O$ of metastable record observables.
\item \textbf{Update capability.} It can condition future local actions on those records.
\item \textbf{Comparison capability.} It can query at least one shared overlap observable
family and detect mismatch syndromes with neighbors.
\item \textbf{Persistence.} The pattern remains identifiable for many local update cycles.
\end{enumerate}

This is an observer pattern criterion, not a consciousness criterion. It is meant to support
simulator definitions of record-bearing internal agents or agent-like subsystems.

\begin{verbatim}
Local registers --> records --> observer pattern --> synchronization

   link qudits / qubits
          |
          v
   gauge-invariant patch observables
          |
          v
   robust pointer variables / records
          |
          v
   persistent update-capable pattern O_ref
          |
          v
   compare shared overlap data and repair mismatch
\end{verbatim}

\section{Interaction, Synchronization, and Repair}

The useful phase-1 object is not a vague ``repair option'' but a finite overlap interface that a
simulator can instantiate on every edge of the patch-overlap graph. For each neighboring patch
pair $(P,Q)$, the present paper should specify the support of the update, the declared shared
observable family, the finite syndrome alphabet, the local correction menu, the acceptance rule,
and the logging contract. What it does \emph{not} yet claim is a theorem that a chosen family
must terminate, be confluent, or yield a continuum-consistent normal form. Those theorem-level
questions remain on the consensus-paper side of the boundary.

The phase-1 default should be explicit syndrome-based overlap repair rather than a fully
implicit dissipative dynamics. Dissipative or record-mediated variants may later be useful, but
they should compile to the same fixed-cutoff interface: read shared data, compare it, choose one
local candidate correction, accept or reject it by a declared contract, then refresh the record
summaries and log the result.

\subsection{Minimal overlap API}

For each neighboring pair $(P,Q)$ with shared collar $B_{PQ}=P\cap Q$, define the one-step
update neighborhood $N_{PQ}$ to be $B_{PQ}$ together with all links, vertices, record registers,
and faces at graph distance at most one from $B_{PQ}$. No single repair step is allowed to act
outside $N_{PQ}$. A phase-1 simulator should instantiate one typed overlap object
\begin{equation}
\mathfrak I_{PQ}
=
\Bigl(
B_{PQ},
N_{PQ},
\mathcal O^{\mathrm{sh}}_{PQ},
\mathsf X_{PQ},
\operatorname{read}_{P\to Q},
\operatorname{read}_{Q\to P},
\operatorname{compare}_{PQ},
\mathcal C^{\mathrm{min}}_{PQ},
\operatorname{accept}_{PQ},
\operatorname{commit}_{PQ}
\Bigr)
\end{equation}
for every edge of the patch-overlap graph. In the explicit builds below, the generic
$\mathcal C^{\mathrm{min}}_{PQ}$ becomes the concrete abelian menu
$\mathcal C^{(2)}_{IJ}$ or the nonabelian menu $\mathcal C^{(S_3)}_{IJ}$.

Its entries should be interpreted literally.
\begin{enumerate}[leftmargin=2em]
\item $B_{PQ}$ is the collar support on which the shared observables live.
\item $N_{PQ}$ is the maximal support on which one local repair move may act.
\item $\mathcal O^{\mathrm{sh}}_{PQ}$ is the declared shared observable family to be
synchronized on this overlap.
\item $\mathsf X_{PQ}$ is the finite alphabet of readout packets exported by each side.
\item $\operatorname{read}_{P\to Q}$ and $\operatorname{read}_{Q\to P}$ are the local readout
channels that turn the current microscopic state into those packets.
\item $\operatorname{compare}_{PQ}$ maps two packets into a finite mismatch syndrome.
\item $\mathcal C^{\mathrm{min}}_{PQ}$ is the finite menu of allowed local corrections.
\item $\operatorname{accept}_{PQ}$ checks whether a candidate correction satisfies the repair
contract.
\item $\operatorname{commit}_{PQ}$ applies the accepted move, refreshes the declared record
summaries, and writes the audit log.
\end{enumerate}

In code, this is the per-overlap declaration table. Once $\mathfrak I_{PQ}$ has been fixed,
different schedules may choose different overlaps or different candidate-order policies, but they
should not silently change the shared observable family, the syndrome alphabet, or the
acceptance rule mid-run.

\subsection{Declared shared observable family and readout packets}

What is synchronized is not the full microscopic quantum state. The synchronization target is a
small declared family of overlap observables,
\begin{equation}
\mathcal O^{\mathrm{sh}}_{PQ}
=
\{\Pi_{\alpha}^{(PQ)}\}_{\alpha\in\mathsf A_{PQ}}
\cup
\{Q_a^{(PQ)}\}_{a=1}^{n_Q}
\cup
\{M_r^{(PQ)}\}_{r=1}^{n_M}
\subseteq
\mathcal A(B_{PQ}).
\end{equation}
For phase 1, every overlap should expose at least one observable from each of the following
three blocks.
\begin{enumerate}[leftmargin=2em]
\item \textbf{Sector block.}
$\{\Pi_{\alpha}^{(PQ)}\}$ is a complete family of collar-sector projectors, with
\begin{equation}
\Pi_{\alpha}^{(PQ)}\Pi_{\beta}^{(PQ)}=\delta_{\alpha\beta}\Pi_{\alpha}^{(PQ)},
\qquad
\sum_{\alpha\in\mathsf A_{PQ}}\Pi_{\alpha}^{(PQ)}=\mathbf 1.
\end{equation}
The inferred sector label $\alpha$ is the first overlap-level datum to compare.
\item \textbf{Boundary observable block.}
Each $Q_a^{(PQ)}$ is a gauge-invariant collar observable with finite value alphabet
$\mathsf V_a$ and declared comparison rule
\begin{equation}
\operatorname{diff}_a:\mathsf V_a\times \mathsf V_a\to \Delta_a.
\end{equation}
For abelian pilots $\operatorname{diff}_a$ may simply be subtraction mod $N$ or a binary
equality test. For nonabelian pilots it may return a class, irrep, or other finite comparison
label fixed in advance.
\item \textbf{Record-summary block.}
Each $M_r^{(PQ)}$ is a read-stable summary variable built from the local record registers and/or
recoverable overlap data, with finite value alphabet $\mathsf W_r$, declared comparison rule
\begin{equation}
\operatorname{diff}_r:\mathsf W_r\times \mathsf W_r\to \Delta_r,
\end{equation}
and freshness threshold $T_r^{\max}$. In the first simulator build these should be chosen from
commuting or numerically stable approximately commuting pointer variables.
\end{enumerate}

The readout channel from the $P$ side of the overlap should return the packet
\begin{equation}
x_{P\to Q}^{(PQ)}
=
\bigl(
\hat g_P,
\hat\alpha_P,
\hat q_{P,1},\dots,\hat q_{P,n_Q},
\hat m_{P,1},\dots,\hat m_{P,n_M},
\hat s_{P,1},\dots,\hat s_{P,n_M}
\bigr)
\in \mathsf X_{PQ},
\end{equation}
and similarly for $x_{Q\to P}^{(PQ)}$. Here:
\begin{enumerate}[leftmargin=2em]
\item $\hat g_P\in\{0,1\}$ is the local admissibility flag on the support $N_{PQ}\cap P$;
\item $\hat\alpha_P\in\mathsf A_{PQ}$ is the inferred collar-sector label;
\item $\hat q_{P,a}\in\mathsf V_a$ are the declared boundary-observable values;
\item $\hat m_{P,r}\in\mathsf W_r$ are the declared record-summary values;
\item $\hat s_{P,r}\in\{0,1\}$ is the stale-bit flag for record summary $r$, with
$\hat s_{P,r}=1$ meaning either that the age of the stored summary exceeds $T_r^{\max}$ or that
the last refresh attempt failed.
\end{enumerate}

A simulator build should treat the packet format as part of the API, not as an informal
comment. Every overlap should declare the alphabets $\mathsf A_{PQ}$, $\mathsf V_a$,
$\mathsf W_r$, the comparison rules $\operatorname{diff}_a$, $\operatorname{diff}_r$, and the
freshness thresholds before the first dynamics run starts.

\subsection{Mismatch syndrome and local acceptance score}

The mismatch syndrome is the finite object returned by the comparison map
\begin{align}
S_{PQ}
&=
\Bigl(
\chi_G,
\chi_{\alpha},
\{\Delta_a^{Q}\}_{a=1}^{n_Q},
\{\Delta_r^{M}\}_{r=1}^{n_M},
\{\sigma_r^{\mathrm{stale}}\}_{r=1}^{n_M}
\Bigr),\\
\chi_G
&=
1-\hat g_P\hat g_Q,\\
\chi_{\alpha}
&=
\mathbf 1[\hat\alpha_P\neq \hat\alpha_Q],\\
\Delta_a^{Q}
&=
\operatorname{diff}_a(\hat q_{P,a},\hat q_{Q,a}),\\
\Delta_r^{M}
&=
\operatorname{diff}_r(\hat m_{P,r},\hat m_{Q,r}),\\
\sigma_r^{\mathrm{stale}}
&=
\hat s_{P,r}\lor \hat s_{Q,r}.
\end{align}
This definition makes the syndrome alphabet completely explicit: one gauge flag, one
sector-mismatch flag, one finite comparison output for each declared boundary observable, one
finite comparison output for each record summary, and one stale-bit flag per record summary.

For local acceptance it is convenient to use an exact lexicographic score rather than a single
weighted scalar:
\begin{equation}
L_{PQ}(S_{PQ})
=
\Bigl(
\chi_G,
\chi_{\alpha},
\sum_{a=1}^{n_Q}\mathbf 1[\Delta_a^{Q}\neq 0],
\sum_{r=1}^{n_M}\mathbf 1[\Delta_r^{M}\neq 0],
\sum_{r=1}^{n_M}\sigma_r^{\mathrm{stale}}
\Bigr),
\end{equation}
ordered lexicographically from left to right. A phase-1 accepted correction should strictly
lower $L_{PQ}$. This hard-codes the priority order
\begin{equation}
\text{gauge admissibility}
\;>\;
\text{sector agreement}
\;>\;
\text{boundary-observable agreement}
\;>\;
\text{record-summary agreement}
\;>\;
\text{record freshness}.
\end{equation}

For logging and later schedule diagnostics, it is still useful to keep a scalar local cost,
\begin{equation}
\Xi_{PQ}
=
\lambda_G\chi_G
+
\lambda_{\alpha}\chi_{\alpha}
+
\sum_{a=1}^{n_Q}\lambda_a\,\mathbf 1[\Delta_a^{Q}\neq 0]
+
\sum_{r=1}^{n_M}\lambda_r\,\mathbf 1[\Delta_r^{M}\neq 0]
+
\lambda_{\mathrm{stale}}\sum_{r=1}^{n_M}\sigma_r^{\mathrm{stale}},
\qquad \lambda_{\bullet}\ge 0,
\end{equation}
and the corresponding global mismatch potential
\begin{equation}
\Phi=\sum_{\langle P,Q\rangle} w_{PQ}\,\Xi_{PQ}.
\end{equation}
The intended use is operational: $L_{PQ}$ determines whether a single move is accepted, while
$\Xi_{PQ}$ and $\Phi$ are diagnostic quantities logged over a run. Nothing in this paper proves
that a chosen schedule must decrease $\Phi$ globally.

\subsection{Repair contract and minimal correction menu}

The minimal correction menu should be intentionally small:
\begin{align}
\mathcal C^{\mathrm{min}}_{PQ}
&=
\mathcal C^{\mathrm{link}}_{PQ}
\cup
\mathcal C^{\mathrm{face}}_{PQ}
\cup
\mathcal C^{\mathrm{rec}}_{PQ}
\cup
\{\mathrm{noop}\},\\
\mathcal C^{\mathrm{link}}_{PQ}
&=
\{C_{e}^{(u)}: e\subset N_{PQ}\cap \partial B_{PQ},\ u\in U_e\},\\
\mathcal C^{\mathrm{face}}_{PQ}
&=
\{C_{f}^{(u)}: f\cap B_{PQ}\neq\varnothing,\ u\in U_f\},\\
\mathcal C^{\mathrm{rec}}_{PQ}
&=
\{C_{r}^{\mathrm{refresh}}: r\text{ supported on }N_{PQ}\}.
\end{align}
Here $U_e$ and $U_f$ are finite move sets fixed by the microscopic model. For the octahedral
$\mathbb Z_2$ pilot, a natural first choice is
\begin{equation}
U_e=\{X_e\},
\qquad
U_f=\Bigl\{\prod_{e\subset\partial f}X_e\Bigr\}.
\end{equation}
For a small nonabelian finite-group pilot, $U_e$ may be left or right multiplication by a
chosen generator set, and $U_f$ the corresponding minimal plaquette moves. The record-refresh
move $C_r^{\mathrm{refresh}}$ is allowed to rewrite only the declared summary variable and its
stale bit from the current post-update overlap data; it is not allowed to invent a new
synchronization target.

A candidate correction $c\in\mathcal C^{\mathrm{min}}_{PQ}$ should satisfy the following repair
contract.
\begin{enumerate}[leftmargin=2em]
\item \textbf{Support contract.}
$\operatorname{supp}(c)\subseteq N_{PQ}$.
\item \textbf{Gauge contract.}
After tentative application of $c$, all Gauss checks on $N_{PQ}$ must pass before the move can
be committed.
\item \textbf{Shared-family contract.}
The move is evaluated only against the predeclared family $\mathcal O^{\mathrm{sh}}_{PQ}$; the
simulator may not enlarge or redefine that family while a run is in progress.
\item \textbf{Monotone local-fit contract.}
If $S_{PQ}^{c}$ is the post-move syndrome, then the move may be accepted only if
\begin{equation}
L_{PQ}(S_{PQ}^{c}) <_{\mathrm{lex}} L_{PQ}(S_{PQ}).
\end{equation}
\item \textbf{Record contract.}
If the move touches a declared record summary, the committed record value must match the
post-update value of the corresponding $M_r^{(PQ)}$, and any stale bit that has actually been
refreshed must be cleared.
\item \textbf{Audit contract.}
Every attempted move must log the pre-packets, post-packets, syndrome, candidate identifier,
support, and accept/reject reason.
\end{enumerate}

A deterministic candidate selector is then
\begin{equation}
c_{\ast}
=
\arg\min_{c\in\mathcal C^{\mathrm{min}}_{PQ}}
\Bigl(
L_{PQ}(S_{PQ}^{c}),
|\operatorname{supp}(c)|,
\operatorname{id}(c)
\Bigr),
\end{equation}
where $\operatorname{id}(c)$ is a fixed total ordering on the candidate list and the minimum is
taken only over candidates satisfying the support and gauge contracts. The move is committed
only if it also satisfies the monotone local-fit contract. The no-op move is accepted only when
$S_{PQ}=0$; otherwise it is logged as an unresolved mismatch, not disguised as success.

\subsection{Local update loop}

Once every overlap object $\mathfrak I_{PQ}$ has been declared, the asynchronous local update
loop should be completely explicit.

\begin{verbatim}
Initialize the microscopic state, the patch cover, the overlap objects I_PQ,
and the record registers.

For t = 1,...,T_max:
    choose one overlap (P,Q) according to schedule nu_t
    x_P <- read_{P->Q}^{(PQ)}
    x_Q <- read_{Q->P}^{(PQ)}
    S   <- compare_{PQ}(x_P, x_Q)

    if S = 0:
        log no-op, Xi_{PQ}, and Phi
        continue

    build all trial candidates c in C_{PQ}^{min} \ {noop}
    discard candidates that violate local Gauss admissibility
    score each surviving candidate by
        ( post-update L_{PQ}, support size, static candidate id )

    if no surviving candidate strictly lowers L_{PQ}:
        log unresolved syndrome and continue

    commit the best candidate
    refresh every touched record summary and stale bit
    reread x'_P, x'_Q and recompute S'_{PQ}, Xi_{PQ}, and Phi
    log the accepted move
\end{verbatim}

Only the schedule $\nu_t$ should vary across the simplest simulation experiments. Reasonable
phase-1 choices are round-robin over overlap edges, uniform random asynchronous sampling, or
priority by current largest $\Xi_{PQ}$. The interface $\mathfrak I_{PQ}$ should remain
unchanged across those schedule experiments.

This loop also makes the simulator/logging boundary concrete. A single log line per attempted
overlap step should at minimum contain the overlap label $(P,Q)$, the two input packets, the
syndrome, the chosen candidate (or explicit failure), the post-update packet values, the local
score $L_{PQ}$ before and after, and the diagnostics $\Xi_{PQ}$ and $\Phi$. Without that log,
later schedule and obstruction analysis is underdetermined.

\subsection{Phase-1 acceptance tests and handoff boundary}

The first simulator build should be required to pass explicit interface-level tests before any
larger interpretation is attached to it.
\begin{enumerate}[leftmargin=2em]
\item \textbf{Schema completeness.}
Every overlap in the patch-overlap graph must ship with declared alphabets $\mathsf A_{PQ}$,
$\mathsf V_a$, $\mathsf W_r$, comparison rules $\operatorname{diff}_a$, $\operatorname{diff}_r$,
freshness thresholds $T_r^{\max}$, and a finite correction menu. Missing declarations are build
failures.
\item \textbf{Zero-syndrome sanity.}
If two overlap packets agree field by field and all stale bits are zero, then
$\operatorname{compare}_{PQ}$ must return $S_{PQ}=0$ and the only legal response is the no-op.
\item \textbf{Single-defect detectability.}
A manually injected mismatch in the sector label, one boundary observable, or one record summary
must flip the corresponding syndrome component and should not generate unrelated syndrome
components except through an explicit gauge violation.
\item \textbf{Gauge-safe commit.}
Every accepted nontrivial correction must leave all Gauss checks on $N_{PQ}$ satisfied.
\item \textbf{Strict local improvement.}
For every accepted correction,
\begin{equation}
L_{PQ}^{\mathrm{after}} <_{\mathrm{lex}} L_{PQ}^{\mathrm{before}}.
\end{equation}
Reject logs must state which repair-contract clause failed.
\item \textbf{Record-refresh fidelity.}
After an accepted refresh move, rereading the overlap must return record-summary values equal to
the current declared $M_r^{(PQ)}$ values, and the corresponding stale bits must clear.
\item \textbf{Schedule comparison on a small cover.}
On the octahedral $\mathbb Z_2$ pilot with a three-patch tree-like cover, rerun the same
initial condition under round-robin, random, and largest-$\Xi$ schedules and compare the final
shared-observable packets. Agreement on that declared family is evidence that the interface is
functioning on that instance; disagreement is a measured failure mode, not a contradiction of
any theorem in this paper.
\item \textbf{Residual-obstruction exposure.}
On a deliberately frustrated cycle, the loop should surface unresolved syndromes or schedule
dependence in the logs rather than silently overwriting record summaries to fake agreement.
\end{enumerate}

This is the intended handoff boundary. The present microphysics paper owns the fixed-cutoff
interface: the overlap object $\mathfrak I_{PQ}$, the declared shared observable family, the
readout packets, the syndrome map, the local correction menu, the acceptance contract, and the
unit/integration tests that a simulator can run. The separate consensus paper should own the
theorem-level questions: termination, local confluence, schedule robustness, normal-form
existence, and the classification of residual holonomy or obstruction classes. Phase 1
therefore supplies the concrete microscopic inputs to the consensus story without claiming that
the global consensus theorems have already been proved here.

\section{What a First Implementable Reference Model Should Contain}

A first actual build should be deliberately small, explicit, and reproducible. The target of the
first round is not a unique UV claim; it is one finite cutoff model with enough declared
structure that two independent teams would compile the same registers, the same layers, the same
logged observables, and the same validation tests.

\subsection{Fixed geometry and patch cover}

Use the octahedral cellulation of $S^2$. Label vertices
\begin{equation}
X^{\pm},\qquad Y^{\pm},\qquad Z^{\pm},
\end{equation}
and faces by sign triples
\begin{equation}
f_{\sigma\tau\upsilon}=(X^{\sigma},Y^{\tau},Z^{\upsilon}),
\qquad \sigma,\tau,\upsilon\in\{+,-\}.
\end{equation}
Thus $|V|=6$, $|E|=12$, and $|F|=8$. Use the fixed vertex order
\begin{equation}
X^+ \prec X^- \prec Y^+ \prec Y^- \prec Z^+ \prec Z^-,
\end{equation}
store one physical link register on each undirected edge, orient that edge from the earlier
endpoint to the later endpoint, and use the same ordered endpoint pair both for register lookup
and for the lexicographic ordering of collar edges in the repair layer. Every holonomy routine
must use that same orientation and invert the group element when the traversal direction is
reversed.

For the octahedral pilot, use the patch cover
\begin{align}
P_X &= \{f_{+\tau\upsilon}:\tau,\upsilon\in\{+,-\}\},\\
P_Y &= \{f_{\sigma+\upsilon}:\sigma,\upsilon\in\{+,-\}\},\\
P_Z &= \{f_{\sigma\tau+}:\sigma,\tau\in\{+,-\}\}.
\end{align}
The pairwise overlaps are
\begin{align}
B_{XY}&=P_X\cap P_Y=\{f_{+++},f_{++-}\},\\
B_{XZ}&=P_X\cap P_Z=\{f_{+++},f_{+-+}\},\\
B_{YZ}&=P_Y\cap P_Z=\{f_{+++},f_{-++}\},
\end{align}
and the triple overlap is the single face $f_{+++}$. This specific cover is the required phase-1
geometry because it already contains pairwise collars, a nontrivial triple-overlap consistency
check, and one face unique to each patch:
\begin{equation}
f_X^{\mathrm{bulk}}=f_{+--},
\qquad
f_Y^{\mathrm{bulk}}=f_{-+-},
\qquad
f_Z^{\mathrm{bulk}}=f_{--+}.
\end{equation}

\subsection{First build: the $\mathbb Z_2$ pilot}

The first complete build should use $G=\mathbb Z_2$ and a stabilizer-friendly discrete update
rule. Put one qubit $q_e$ on each of the $12$ stored links, with
\begin{equation}
|0\rangle_e,\ |1\rangle_e
\end{equation}
the $Z_e$-eigenbasis. Use
\begin{equation}
A_v=\prod_{e\ni v}X_e,
\qquad
P_v=\frac{1+A_v}{2},
\qquad
B_f=\prod_{e\subset\partial f}Z_e.
\end{equation}
For the phase-1 pilot, the required record layer is not one qubit per patch but three qubits per
patch,
\begin{equation}
R_X=(r_X^{\mathrm{bulk}},r_X^{(Y)},r_X^{(Z)}),\quad
R_Y=(r_Y^{\mathrm{bulk}},r_Y^{(X)},r_Y^{(Z)}),\quad
R_Z=(r_Z^{\mathrm{bulk}},r_Z^{(X)},r_Z^{(Y)}),
\end{equation}
for a total of $9$ record qubits. Their declared meanings are
\begin{align}
Z_{r_I^{\mathrm{bulk}}} &\leftrightarrow B_{f_I^{\mathrm{bulk}}},\\
Z_{r_I^{(J)}} &\leftrightarrow C_{IJ}:=\prod_{e\in \partial B_{IJ}} Z_e,
\end{align}
where $C_{IJ}$ is the four-edge collar-loop parity around the two-face overlap $B_{IJ}$.

The required initial state is any explicitly prepared state in the Gauss-invariant subspace. The
default debug initialization is
\begin{equation}
|\Omega_0\rangle \propto
\Bigl(\prod_{v\in V}P_v\Bigr)
\Bigl(\prod_{f\in F}\frac{1+B_f}{2}\Bigr)
|+\rangle^{\otimes 12}\otimes |0\rangle_{\mathrm{rec}}^{\otimes 9},
\end{equation}
followed by one record-write layer so that all record qubits contain the declared summaries of
the initialized link state.

\subsection{Required local layers for the $\mathbb Z_2$ pilot}

The first build should be a layered discrete-time simulator. One full cycle applies the layers in
the order
\begin{equation}
L_{\mathrm{bulk}}
\longrightarrow
L_{\mathrm{check}}
\longrightarrow
L_{\mathrm{write}}
\longrightarrow
L_{\mathrm{compare}}
\longrightarrow
L_{\mathrm{repair}}
\longrightarrow
L_{\mathrm{verify}}.
\end{equation}
The layers are fixed as follows.
\begin{enumerate}[leftmargin=2em]
\item \textbf{Bulk face-update layer $L_{\mathrm{bulk}}$.}
On a chosen active face set $F_{\mathrm{act}}(t)$, apply either the identity or the face move
\begin{equation}
F_f:=\prod_{e\subset\partial f}X_e.
\end{equation}
The default schedule is a deterministic scan over the eight faces with one Bernoulli face-move
attempt per face and default attempt rate $p_{\mathrm{face}}=0.1$.
\item \textbf{Gauge-check layer $L_{\mathrm{check}}$.}
Measure or compute all six star values
\begin{equation}
a_v(t)\in\{+1,-1\},
\qquad a_v(t)=\langle A_v\rangle_t
\end{equation}
in the exact backend, or their sampled eigenvalues in the shot backend. Because the allowed
phase-1 bulk and repair moves all commute with every $A_v$, any nonzero star syndrome in the
pilot is treated as an implementation fault, not as physical dynamics.
\item \textbf{Record-write layer $L_{\mathrm{write}}$.}
Overwrite the record qubits with the declared summaries
\begin{align}
r_I^{\mathrm{bulk}} &\leftarrow B_{f_I^{\mathrm{bulk}}},\\
r_I^{(J)} &\leftarrow C_{IJ},
\end{align}
using parity-compute and uncompute subcircuits or their exact-state equivalents. The required
mode for the first build is overwrite, not additive update, so that stale-record diagnostics are
unambiguous.
\item \textbf{Overlap-compare layer $L_{\mathrm{compare}}$.}
Choose one active overlap from
\begin{equation}
(XY),(XZ),(YZ)
\end{equation}
by cyclic schedule or by a logged pseudorandom seed. Measure the live shared observables on that
overlap:
\begin{equation}
c_{IJ}=C_{IJ},\qquad
b_{IJ}^{+}=B_{f_{IJ}^{+}},\qquad
b_{IJ}^{-}=B_{f_{IJ}^{-}},
\end{equation}
together with the two cached record values
\begin{equation}
m_I^{(J)}=Z_{r_I^{(J)}},\qquad
m_J^{(I)}=Z_{r_J^{(I)}}.
\end{equation}
For the named overlaps the face labels are
\begin{align}
(f_{XY}^{+},f_{XY}^{-})&=(f_{+++},f_{++-}),\\
(f_{XZ}^{+},f_{XZ}^{-})&=(f_{+++},f_{+-+}),\\
(f_{YZ}^{+},f_{YZ}^{-})&=(f_{+++},f_{-++}).
\end{align}
\item \textbf{Repair layer $L_{\mathrm{repair}}$.}
Evaluate the explicit candidate menu
\begin{equation}
\mathcal C_{IJ}^{(2)}
=
\{\mathrm{Id},\ \mathrm{Refresh}_I,\ \mathrm{Refresh}_J,\ \mathrm{Refresh}_{IJ}\}
\cup
\{X_e:e\in \partial B_{IJ}\}
\cup
\{F_f:f\subset B_{IJ}\}.
\end{equation}
For each candidate $C\in\mathcal C_{IJ}^{(2)}$, compute the post-correction local cost
\begin{equation}
\Xi_{IJ}^{(2)}(C)
=
\lambda_c\,\mathbf 1[m_I^{(J)}\neq c_{IJ}]_C
+
\lambda_c\,\mathbf 1[m_J^{(I)}\neq c_{IJ}]_C
+
\lambda_r\,\mathbf 1[m_I^{(J)}\neq m_J^{(I)}]_C,
\end{equation}
with default weights
\begin{equation}
\lambda_c=2,
\qquad \lambda_r=1.
\end{equation}
Choose the unique minimizer if one exists. If there is a tie, use the fixed tie-break order
\begin{equation}
\mathrm{Refresh}_I \prec \mathrm{Refresh}_J \prec \mathrm{Refresh}_{IJ}
\prec X_{e_1}\prec X_{e_2}\prec X_{e_3}\prec X_{e_4}
\prec F_{f_{IJ}^{+}}\prec F_{f_{IJ}^{-}}\prec \mathrm{Id},
\end{equation}
where $e_1,\dots,e_4$ is the lexicographic ordering of the four collar edges induced by the
stored endpoint order above. This tie-break order must be logged and fixed across all runs.
\item \textbf{Verification layer $L_{\mathrm{verify}}$.}
Recompute $a_v$, the active overlap data, and the global mismatch potential
\begin{equation}
\Phi(t)=\Xi_{XY}^{(2)}(t)+\Xi_{XZ}^{(2)}(t)+\Xi_{YZ}^{(2)}(t).
\end{equation}
A cycle is accepted only after the post-repair measurements have been written to the log.
\end{enumerate}

The designated observer proxy for the pilot is patch $P_X$: its controller is required to choose
the active correction on $B_{XY}$ or $B_{XZ}$ using only the locally available record bits in
$R_X$ plus the currently exposed overlap data. That is enough to instantiate the paper's
operational observer criterion at fixed cutoff without making stronger claims.

\subsection{First nonabelian upgrade: the $S_3$ build}

The first upgrade should keep the same octahedral graph, the same patch cover, the same overlap
schedule, and the same log schema. Only the link and record alphabets change. Each link register
becomes a six-state $S_3$ register encoded into three qubits with the fixed lookup table
\begin{align}
000&\leftrightarrow e, &
001&\leftrightarrow (12), &
010&\leftrightarrow (13),\\
011&\leftrightarrow (23), &
100&\leftrightarrow (123), &
101&\leftrightarrow (132),
\end{align}
while $110$ and $111$ are illegal code words. Reverse edge traversal uses the inverse group
element, implemented by the fixed inversion lookup
\begin{equation}
e^{-1}=e,\qquad (12)^{-1}=(12),\qquad (13)^{-1}=(13),\qquad (23)^{-1}=(23),\qquad
(123)^{-1}=(132).
\end{equation}

At each vertex, the required Gauss projector is the exact finite-group projector
\begin{equation}
P_v^{S_3}
=
\frac{1}{6}\sum_{h\in S_3} U_v(h),
\end{equation}
with $U_v(h)$ acting by left or right multiplication on the incident oriented links according to
their incoming or outgoing orientation. For each face $f$, compute the ordered face holonomy
\begin{equation}
h_f=\prod_{e\subset \partial f} g_e^{\epsilon(e,f)},
\end{equation}
and for each overlap collar compute
\begin{equation}
h_{IJ}^{\mathrm{collar}}=\prod_{e\in \partial B_{IJ}} g_e^{\epsilon(e,IJ)}.
\end{equation}
The required record layer for the first $S_3$ build is binary-coded rather than one-hot. Each
patch $P_I$ carries
\begin{equation}
R_I^{S_3}=
\bigl(c_I^{\mathrm{bulk}},c_I^{(J_1)},c_I^{(J_2)},v_I\bigr),
\end{equation}
where each $c$-register is a two-qubit class label with encoding
\begin{equation}
00\leftrightarrow [e],\qquad
01\leftrightarrow [\tau],\qquad
10\leftrightarrow [\sigma],\qquad
11\ \text{illegal},
\end{equation}
and $v_I$ is a one-bit valid-write flag set to $1$ only when the current write produced a legal
class label. Here $[\tau]$ is the transposition class and $[\sigma]$ the three-cycle class.

The local layer structure is the same as in the $\mathbb Z_2$ pilot, but the candidate correction
set is now
\begin{align}
\mathcal C_{IJ}^{(S_3)}
=&
\{\mathrm{Id},\ \mathrm{Refresh}_I,\ \mathrm{Refresh}_J,\ \mathrm{Refresh}_{IJ}\}\notag\\
&\cup
\{L_e(s),L_e(r),L_e(r^{-1}),R_e(s),R_e(r),R_e(r^{-1}) : e\in \partial B_{IJ}\}\notag\\
&\cup
\{M_f(s),M_f(r),M_f(r^{-1}) : f\subset B_{IJ}\},
\end{align}
where $s=(12)$, $r=(123)$, $L_e(g)$ and $R_e(g)$ left- or right-multiply one collar edge by a
generator, and $M_f(g)$ multiplies one chosen face boundary by the same group element in the
oriented order. The default nonabelian local cost is
\begin{equation}
\Xi_{IJ}^{(S_3)}(C)
=
\lambda_c\,\mathbf 1[c_I^{(J)}\neq \mathrm{cl}(h_{IJ}^{\mathrm{collar}})]_C
+
\lambda_c\,\mathbf 1[c_J^{(I)}\neq \mathrm{cl}(h_{IJ}^{\mathrm{collar}})]_C
+
\lambda_r\,\mathbf 1[c_I^{(J)}\neq c_J^{(I)}]_C
+
\lambda_u\,N_{\mathrm{illegal}}(C),
\end{equation}
with defaults
\begin{equation}
\lambda_c=2,\qquad \lambda_r=1,\qquad \lambda_u=10.
\end{equation}
In addition to conjugacy-class data, the first $S_3$ upgrade must also log the three collar
irrep projectors
\begin{equation}
\Pi_{\rho}^{(IJ)}
=
\frac{d_{\rho}}{6}\sum_{g\in S_3}\chi_{\rho}(g^{-1})\,U_g^{(IJ)},
\qquad
\rho\in\{\mathbf 1,\mathbf 1',\mathbf 2\},
\end{equation}
because the whole point of the upgrade is to check that the overlap interface still behaves well
when the shared sector carries genuinely nonabelian labels.

\subsection{Required deliverables of the first build package}

A phase-1 implementation team should treat the following outputs as mandatory rather than
optional.
\begin{enumerate}[leftmargin=2em]
\item A machine-readable geometry file containing the vertex list, oriented edge list, face list,
patch memberships, and overlap memberships for the octahedron.
\item A register file declaring every link register, every record register, and every illegal-code
word.
\item A layer file specifying $L_{\mathrm{bulk}},L_{\mathrm{check}},L_{\mathrm{write}},
L_{\mathrm{compare}},L_{\mathrm{repair}},L_{\mathrm{verify}}$ in the exact order used above.
\item A correction file containing the candidate menus $\mathcal C_{IJ}^{(2)}$ and
$\mathcal C_{IJ}^{(S_3)}$, the weights $(\lambda_c,\lambda_r,\lambda_u)$, and the tie-break
ordering.
\item A log schema that writes raw register states, derived observables, selected corrections,
and pass/fail statistics for every cycle.
\end{enumerate}

This remains a phase-1 architecture note. A successful pilot would show that one fixed-cutoff
screen model can instantiate the declared objects and measurements; it would not identify a
unique UV completion or close any continuum claim.

\section{Simulator Validation Package}

A useful simulator package has to be able to fail. In phase 1 the validation target is not
continuum closure. It is only the fixed-cutoff architecture claimed in this note: gauge-admissible
local layers, measurable overlap observables, persistent records, and a repair loop whose
behavior can be diagnosed. The validation package should therefore be distributed with the first
build, not added later as an informal plotting appendix.

\subsection{Required benchmark families}

Every implementation must run the same named benchmark suite. The minimum required suite is:
\begin{enumerate}[leftmargin=2em]
\item \texttt{Z2-0-clean}: initialize the $\mathbb Z_2$ pilot in $|\Omega_0\rangle$, run $12$
cycles, and inject no defect.
\item \texttt{Z2-1-stale-record}: after the first $L_{\mathrm{write}}$, flip $r_X^{(Y)}$ and
leave the link state untouched; run $12$ cycles.
\item \texttt{Z2-2-collar-flip}: before the first $L_{\mathrm{compare}}$ on $B_{XY}$, apply
$X_e$ on the lexicographically first collar edge of $B_{XY}$; run $12$ cycles.
\item \texttt{Z2-3-double-overlap}: inject one collar-edge flip on $B_{XY}$ and one on $B_{YZ}$
before the first compare layer; run $15$ cycles.
\item \texttt{Z2-4-frustrated-cycle}: after every write layer during cycles $0$ through $4$,
force $r_X^{(Y)}$ to the opposite of the live collar bit on $B_{XY}$; run $15$ cycles to test
persistent residual mismatch.
\item \texttt{S3-0-clean}: initialize the $S_3$ build in the declared identity sector, run $15$
cycles, and inject no defect.
\item \texttt{S3-1-transposition}: before the first compare layer on $B_{XY}$, left-multiply the
lexicographically first collar edge by $(12)$; run $20$ cycles.
\item \texttt{S3-2-threecycle}: before the first compare layer on $B_{XY}$, left-multiply the
same edge by $(123)$; run $20$ cycles.
\item \texttt{S3-3-mixed-class-cycle}: inject a transposition-class defect on $B_{XY}$ and a
three-cycle defect on $B_{YZ}$ before the first compare layer; run $20$ cycles.
\end{enumerate}

For each benchmark, run all six sweep orders of the three overlaps $(XY,XZ,YZ)$. The analytic
backend must run each order once exactly. The sampled backend must run each order with $32$
independent seeds and $1024$ shots per seed. A reported phase-1 pass is not valid unless it
includes all named benchmarks and all six schedule orders.

\subsection{Required run manifest and per-cycle log}

Every run must write two artifacts.

\paragraph{Run manifest.}
The manifest is one record with the fields \texttt{build\_id}, \texttt{group},
\texttt{backend}, \texttt{encoding}, \texttt{graph\_id},
\texttt{patch\_cover\_id}, \texttt{benchmark\_id}, \texttt{schedule\_id},
\texttt{seed}, \texttt{shots}, \texttt{cycles},
\texttt{candidate\_menu\_hash}, and \texttt{tie\_break\_hash}.
The hashes are required so that two runs can be compared without guessing whether the correction
table changed.

\paragraph{Cycle table.}
Each cycle must write one row with the fields \texttt{cycle},
\texttt{active\_overlap}, \texttt{link\_state\_pre},
\texttt{record\_state\_pre}, \texttt{a\_v\_pre[6]},
\texttt{face\_obs\_pre[8]}, \texttt{collar\_obs\_pre[3]},
\texttt{candidate\_costs}, \texttt{chosen\_correction}, \texttt{accepted},
\texttt{tie\_break\_rank}, \texttt{link\_state\_post},
\texttt{record\_state\_post}, \texttt{a\_v\_post[6]},
\texttt{face\_obs\_post[8]}, \texttt{collar\_obs\_post[3]},
\texttt{Phi\_post}, \texttt{g\_post}, and \texttt{u\_bad\_post}.
For the $\mathbb Z_2$ pilot, \texttt{link\_state} is a $12$-bit string in stored-edge order and
\texttt{face\_obs} is the ordered list of eight plaquette bits $B_f$. For the $S_3$ build,
\texttt{link\_state} must contain both the raw three-bit code word and the decoded group symbol
on every edge, while \texttt{face\_obs} must contain the eight decoded face holonomies together
with their conjugacy classes. For the $S_3$ build, \texttt{collar\_obs} must additionally
contain the three irrep-probability triples
\begin{equation}
\bigl(p_{\mathbf 1}^{(IJ)},p_{\mathbf 1'}^{(IJ)},p_{\mathbf 2}^{(IJ)}\bigr).
\end{equation}

The cycle table is not optional. A run that stores only aggregate plots is a failed validation
artifact because it cannot support repair-level diagnosis.

\subsection{Derived measured quantities}

Let
\begin{equation}
c_{IJ}(t)=
\begin{cases}
C_{IJ}(t), & \mathbb Z_2,\\
\mathrm{cl}(h_{IJ}^{\mathrm{collar}}(t)), & S_3,
\end{cases}
\end{equation}
and let $m_I^{(J)}(t)$ be the corresponding cached overlap record. From the manifest and cycle
table, compute at least the following quantities. Here $\Xi_{IJ}(t)$ denotes the
build-appropriate local mismatch cost on overlap $B_{IJ}$, namely
$\Xi_{IJ}^{(2)}(t)$ in the $\mathbb Z_2$ pilot and $\Xi_{IJ}^{(S_3)}(t)$ in the
$S_3$ build:
\begin{align}
g(t) &= \sum_{v\in V}\mathbf 1[a_v(t)=-1],\\
u_{\mathrm{bad}}(t) &= N_{\mathrm{illegal\ link\ codes}}(t)+N_{\mathrm{illegal\ record\ codes}}(t),\\
\Phi(t) &= \sum_{IJ\in\{XY,XZ,YZ\}}\Xi_{IJ}(t),\\
G_{\Phi} &= \frac{1}{N_{\mathrm{acc}}}\sum_{t\in \mathrm{accepted}}\bigl(\Phi(t-1)-\Phi(t)\bigr),\\
C_r(k) &= \frac{1}{N_{\mathrm{rec}}}\sum_r \Pr[r(t_0)=r(t_0+k)],\\
Q_{\mathrm{readback}} &= \max_{IJ}\|\mu_{IJ}^{\mathrm{read}}-\mu_{IJ}^{\mathrm{skip}}\|_1,\\
D_{\mathrm{sched}} &= \max_{s,s'}\sum_{IJ}\|\mu_{IJ}^{(s)}-\mu_{IJ}^{(s')}\|_1,\\
M_{\mathrm{class}} &= \frac{1}{N_{\mathrm{meas}}}\sum \mathbf 1[\widehat{\mathrm{cl}}(h_{IJ}^{\mathrm{collar}})\neq \mathrm{cl}(h_{IJ}^{\mathrm{collar}})],\\
N_{\mathrm{irrep}} &= \max_{IJ,t}\left|p_{\mathbf 1}^{(IJ)}(t)+p_{\mathbf 1'}^{(IJ)}(t)+p_{\mathbf 2}^{(IJ)}(t)-1\right|.
\end{align}
Here $G_{\Phi}$ is the mean accepted improvement in the mismatch potential, $C_r(k)$ is the
record-retention curve, $Q_{\mathrm{readback}}$ measures back-action from an inserted record-read
layer, $D_{\mathrm{sched}}$ measures schedule dependence on the declared shared observable
family, $M_{\mathrm{class}}$ measures nonabelian class-label errors, and $N_{\mathrm{irrep}}$
checks whether the logged irrep probabilities even define a normalized decomposition.

In addition, every run must report the convergence time $N_{\mathrm{conv}}$, defined for
removable-defect benchmarks as the first cycle after which $\Phi=0$ for the rest of the run, and
the residual-obstruction plateau time $N_{\mathrm{plat}}$, defined for frustrated benchmarks as
the first cycle after which both the residual label and the residual value of $\Phi$ stop
changing.

\subsection{Phase-1 pass and fail gates}

The validation package should use explicit gates rather than qualitative language.
\begin{enumerate}[leftmargin=2em]
\item \textbf{Kinematic gate.}
In the analytic backend, every post-verify cycle must satisfy
\begin{equation}
g(t)=0,
\qquad
u_{\mathrm{bad}}(t)=0.
\end{equation}
In the sampled backend, the required bounds are
\begin{equation}
\max_t \frac{g(t)}{|V|}\le 10^{-3},
\qquad
\max_t \frac{u_{\mathrm{bad}}(t)}{N_{\mathrm{reg}}}\le 10^{-4},
\end{equation}
where $N_{\mathrm{reg}}$ is the number of encoded link and record code words checked for
legality in that build.
\item \textbf{Record gate.}
On \texttt{Z2-0-clean}, \texttt{Z2-1-stale-record}, and \texttt{S3-0-clean}, the required record
metrics are
\begin{equation}
C_r(5)\ge 0.80,
\qquad
Q_{\mathrm{readback}}\le 0.05.
\end{equation}
If these fail while the kinematic gate passes, the architecture has local observables but not a
usable record layer.
\item \textbf{$\mathbb Z_2$ synchronization gate.}
On \texttt{Z2-1-stale-record}, \texttt{Z2-2-collar-flip}, and \texttt{Z2-3-double-overlap}, the
analytic backend must reach $\Phi=0$ within $12$ repair activations for all six schedules. The
sampled backend must do so within $15$ repair activations in at least $95\%$ of runs. In both
backends, one also requires
\begin{equation}
G_{\Phi}>0,
\qquad
D_{\mathrm{sched}}\le
\begin{cases}
10^{-8}, & \text{analytic},\\
0.05, & \text{sampled}.
\end{cases}
\end{equation}
\item \textbf{Residual-obstruction gate.}
On \texttt{Z2-4-frustrated-cycle}, the build is not supposed to fake full agreement. The
required behavior is
\begin{equation}
g(t)=0\ \text{for all post-verify cycles},
\qquad
N_{\mathrm{plat}}\le 8,
\end{equation}
and the residual label must remain unchanged for the final $10$ cycles. A run that oscillates
indefinitely, or that reports $\Phi=0$ while the forced stale source is still present, fails
this gate.
\item \textbf{$S_3$ nonabelian gate.}
On \texttt{S3-1-transposition}, \texttt{S3-2-threecycle}, and
\texttt{S3-3-mixed-class-cycle}, the required label-fidelity bounds are
\begin{equation}
M_{\mathrm{class}}\le
\begin{cases}
0, & \text{analytic},\\
0.05, & \text{sampled},
\end{cases}
\qquad
N_{\mathrm{irrep}}\le
\begin{cases}
10^{-10}, & \text{analytic},\\
0.02, & \text{sampled}.
\end{cases}
\end{equation}
For removable class defects, the analytic backend must reach matching collar classes within $20$
repair activations for all six schedules; the sampled backend must do so within $24$ repair
activations in at least $90\%$ of runs. The sampled nonabelian schedule bound is
\begin{equation}
D_{\mathrm{sched}}\le 0.10.
\end{equation}
\item \textbf{Harness gate.}
Every negative control listed below must fail at least one of the gates above. If the negative
controls do not fail, the validation harness is too weak to trust.
\end{enumerate}

The status labels should be reported explicitly. Passing the kinematic, record, and
$\mathbb Z_2$ synchronization gates yields an \texttt{abelian-pilot-pass}. Passing those plus
the nonabelian gate yields a \texttt{nonabelian-upgrade-pass}. A full
\texttt{phase1-architecture-pass} requires both of those together with a passed harness gate.

\subsection{Mandatory negative controls}

The first validation package should contain the following explicit negative controls.
\begin{enumerate}[leftmargin=2em]
\item \texttt{NC-1-no-refresh}: remove $\mathrm{Refresh}_I$, $\mathrm{Refresh}_J$, and
$\mathrm{Refresh}_{IJ}$ from the candidate menu. Expected failure:
\texttt{Z2-1-stale-record} does not converge and $C_r(5)$ drops.
\item \texttt{NC-2-wrong-orientation}: reverse one stored edge orientation in the $S_3$ holonomy
routine without updating the inverse lookup. Expected failure: $M_{\mathrm{class}}$ and
$D_{\mathrm{sched}}$ jump while the $\mathbb Z_2$ pilot may still look healthy.
\item \texttt{NC-3-no-illegal-penalty}: set $\lambda_u=0$ in the $S_3$ build. Expected failure:
$u_{\mathrm{bad}}$ and $N_{\mathrm{irrep}}$ become nonzero.
\item \texttt{NC-4-no-write}: skip the record-write layer. Expected failure: $C_r(5)$ and
$Q_{\mathrm{readback}}$ fail even on clean runs.
\end{enumerate}

\subsection{Failure diagnostics}

Every failed gate must be accompanied by a minimal evidence packet containing the benchmark
identifier, schedule identifier, cycle index, pre- and post-correction register snapshots, the
full candidate-cost table, the chosen correction, and the tie-break rank. The default diagnostic
codes are:
\begin{enumerate}[leftmargin=2em]
\item \textbf{F1 gauge drift.} Signature: $g(t)>0$ after $L_{\mathrm{verify}}$. Likely causes:
a correction move that does not preserve the Gauss sector, a wrong stored-edge ordering, or a
verify layer that is reading the pre-repair state.
\item \textbf{F2 code leakage.} Signature: $u_{\mathrm{bad}}(t)>0$. Likely causes: an
incomplete $S_3$ decode table, a missing invalid-word penalty, or a record class encoder that
writes the illegal word $11$.
\item \textbf{F3 record volatility.} Signature: low $C_r(5)$ with acceptable $g(t)$. Likely
causes: the write layer is not pure overwrite, the record circuit reaches beyond the declared
local support, or the readout channel is feeding back into nonshared plaquettes.
\item \textbf{F4 non-monotone repair.} Signature: $G_{\Phi}\le 0$ or accepted moves that
increase $\Phi$. Likely causes: candidate costs evaluated on stale pre-repair data, a tie-break
order that masks a lower-cost move, or a correction menu that is too small for the declared
syndrome family.
\item \textbf{F5 schedule split.} Signature: small $g(t)$ and reasonable records but large
$D_{\mathrm{sched}}$. Likely causes: missing triple-overlap information, nonconfluent local
tie-break rules, or an overlap observable family that is larger than the repair menu can
support.
\item \textbf{F6 nonabelian label failure.} Signature: the $\mathbb Z_2$ pilot passes but
$M_{\mathrm{class}}$ or $N_{\mathrm{irrep}}$ fails. Likely causes: an orientation mistake in
face or collar holonomies, an incorrect multiplication table, or a bad character-projector
implementation.
\item \textbf{F7 false obstruction or missed obstruction.} Signature:
\texttt{Z2-4-frustrated-cycle} either converges spuriously to $\Phi=0$ or never reaches a stable
plateau. Likely causes: the residual-label logic is not implemented, or the compare layer is
failing to distinguish removable mismatch from a persistent external fault source.
\end{enumerate}

\subsection{Reading rule for the validation package}

Even a full \texttt{phase1-architecture-pass} would not prove the OPH continuum branch, would
not prove the final repair law, and would not identify a unique UV completion. It would only
show that the fixed-cutoff architecture claimed in this note has been instantiated concretely
enough to measure and debug. Conversely, failure is informative. A failed gate tells the
implementation team which layer is not doing the job the paper assigns to it, and the required
logs above tell them where to start looking.

\section{Dependence on Existing OPH Results and Open Bridges}

This note is intentionally one-sided in its logical role. It imports OPH structure that already
exists elsewhere in the program, adds one concrete fixed-cutoff reference architecture on top of
that structure, and leaves the scaling-limit and theorem-level bridges explicit rather than
implicit. The intended claim boundary is therefore
\begin{equation}
\text{imported OPH framework}
\;\longrightarrow\;
\text{phase-1 reference architecture}
\;\longrightarrow\;
\text{later bridge theorems}.
\end{equation}

\subsection{Imported OPH structure used here without re-derivation}

The following ingredients are reused as background OPH structure or already-existing OPH
results/conventions. Their status is inherited from the papers that establish them; this note
does not upgrade or re-prove them.
\begin{enumerate}[leftmargin=2em]
\item The screen-net viewpoint on a horizon screen $S^2$, together with patch algebras and the
language of overlap consistency.
\item The collar and edge-center completion package: shared collar algebras, center/sector
projectors $\Pi_{\alpha}$, and the idea that overlap-accessible data live in a boundary-visible
sector decomposition.
\item The recoverable generalized entropy and Markov/recovery vocabulary in which conditional
mutual information, collar recovery, and edge data are meaningful organizing notions.
\item The local MaxEnt plus refinement-stability branch as the current OPH dynamical background
against which regulated architectures are later judged.
\item The OPH observer-access tuple $(P_O,\mathcal A(P_O),\rho_O,R_O)$, understood as an
existing access model for a record-bearing internal subsystem.
\end{enumerate}

\subsection{New claims made only at the architecture level in this note}

Relative to that imported framework, the genuinely new content of this side note is narrower and
strictly phase-1.
\begin{enumerate}[leftmargin=2em]
\item It selects one finite gauge-register or quantum-link style reference architecture for a
cellulated screen, explicitly as a working model rather than as the unique OPH UV completion.
\item It gives an explicit fixed-cutoff dictionary from local registers to patch observables,
overlap observables, collar-sector readouts, record variables, observer interfaces, and mismatch
syndromes.
\item It specifies one phase-1 repair interface
\begin{equation}
\mathcal R_{PQ}=\mathcal W_{PQ}\circ \mathcal C_{PQ}\circ \mathcal M_{PQ}\circ \mathcal P_{PQ},
\end{equation}
together with a finite local correction-menu template and an asynchronous update loop.
\item It proposes a concrete simulator ladder: an octahedral $\mathbb Z_2$ pilot first and an
$S_3$ upgrade second.
\item It states a validation package that can fail, built from gauge-violation rate, overlap
mismatch, record persistence, and schedule-divergence diagnostics.
\end{enumerate}

These are architecture claims and simulator-interface claims. They are not continuum theorems,
not measurement theorems, not continuation theorems, and not quantitative Standard Model
closures.

\subsection{Boundary markers between imported structure and new architecture content}

Three boundaries matter for reading the paper correctly.
\begin{enumerate}[leftmargin=2em]
\item The existence of collar-sector observables is imported OPH regulator language; the new
claim here is only that the chosen finite-register architecture can be organized so that a
declared family of fixed-cutoff observables plays that interface role.
\item The tuple $(P_O,\mathcal A(P_O),\rho_O,R_O)$ is imported OPH observer language; the new
content here is the added update interface $\mathcal U_O$ and the operational phase-1 criterion
for when a persistent pattern counts as an observer-like subsystem in simulation.
\item The need for overlap repair or synchronization is imported from the broader OPH consistency
program; the new content here is only the explicit microscopic syntax of syndromes, local costs,
and correction menus on a finite screen.
\end{enumerate}

So the paper does \emph{not} claim ``OPH microphysics solved.'' It claims only that one concrete
microphysical build surface has been made explicit enough that later OPH tasks can either prove
things about it, measure things on it, or reject it.

\subsection{Open bridges left deliberately open}

The main unfinished bridges are now supposed to be visible rather than hidden.
\begin{enumerate}[leftmargin=2em]
\item \textbf{Microphysics to quantitative sector weights.}
The note does not derive the edge-sector probability law used downstream in the compact
quantitative branch. It only supplies finite models on which candidate sector-weight laws can be
defined and tested.
\item \textbf{Microphysics to regulated OPH embedding.}
The note does not prove that the reference patch net is the regulated OPH patch net. It only
gives a finite architecture whose objects are shaped so that such an embedding can later be
formulated precisely.
\item \textbf{Microphysics to repair theorems.}
The note does not prove termination, local confluence, gauge covariance, or schedule-independence
of the repair law. It only specifies the microscopic observables and local update primitives on
which those proofs would act.
\item \textbf{Microphysics to measurement theory.}
The note does not prove that the record layer becomes the central or approximately central record
algebra needed for a Born-rule package. It only supplies candidate pointer observables,
overlap summaries, and stability tests.
\item \textbf{Microphysics to observer continuation.}
The note does not prove a restoration map, identity criterion, or backup theorem. It only
isolates the checkpoint-like data and collar interfaces that such a theorem package would have to
use.
\item \textbf{Microphysics to continuum/gravity.}
The note does not derive the OPH continuum, geometric modular branch, or gravity branch from the
present reference architecture. Compatibility is the working hope; derivation remains later work.
\end{enumerate}

That is the intended asymmetry of this section: imported OPH structure is treated as imported;
the finite reference architecture is treated as new phase-1 scaffolding; and every bridge from
that scaffolding to later theorem packages remains explicitly open.

\section{Relation Back to the Existing Completion Tree}

This side paper sits \emph{upstream} of several completion-tree leaves, but it does not by
itself close any of them. Its role is to make the finite-cutoff objects concrete enough that
later task-specific papers can stop talking in placeholders and start proving or testing
statements about explicit regulated data.

For the current microphysics completion workflow, the task-specific subsections below are also
the temporary audit surface for the mapped leaf tasks. Audit happens here first, and the compact,
consensus, and observer papers absorb the approved merge only after that audit is complete.

\subsection{What it feeds now for \texttt{papers.compact.e.24}}\label{sec:microphysics-feeds-e24}

The task \texttt{papers.compact.e.24} asks for the edge heat-kernel / Casimir law to be proved
from OPH microphysics. This note feeds that task now in the following limited sense.
\begin{enumerate}[leftmargin=2em]
\item It supplies concrete finite screen microphysics candidates---first $\mathbb Z_2$, then
$S_3$, and later quantum-link style extensions---on which collar sectors are explicit observables
rather than abstract labels.
\item It supplies declared overlap observables $\Pi_{\alpha}^{(PQ)}$ and hence a concrete place
to define sector frequencies, transition data, and finite-cutoff sector weights.
\item It supplies a simulator protocol that can log those sector statistics while simultaneously
checking gauge admissibility, record stability, and synchronization behavior.
\end{enumerate}
What remains open for \texttt{papers.compact.e.24} is the main theorem burden: this note does
\emph{not} derive the actual heat-kernel / Casimir weight law, does \emph{not} derive the
transmutation parameter or scaling law behind it, does \emph{not} show refinement-stable
universality, and does \emph{not} establish that this reference architecture is the correct or
unique OPH microphysics. It provides the regulated measurement surface, not the quantitative
derivation.

\subsection{What it feeds now for \texttt{papers.reality.a.01}}\label{sec:microphysics-feeds-a01}

The task \texttt{papers.reality.a.01} asks for the finite patch-net model of the consensus paper
to be embedded into the OPH regulator. This note feeds that task now by making the finite patch
net explicit rather than schematic.
\begin{enumerate}[leftmargin=2em]
\item It identifies the regulated patch net: a cellulated screen with explicit local Hilbert
spaces, physical subspace projectors, patch regions, overlap collars, and shared observable
families.
\item It makes the interface alphabet concrete through
\begin{equation}
\mathcal O^{\mathrm{sh}}_{PQ}
=
\{\Pi_{\alpha}^{(PQ)}\}_{\alpha}
\cup
\{Q_a^{(PQ)}\}_a
\cup
\{M_r^{(PQ)}\}_r,
\end{equation}
together with the mismatch syndrome $S_{PQ}$, local cost $\Xi_{PQ}$, and aggregate potential
$\Phi$.
\item It supplies a local repair syntax
\begin{equation}
\mathcal R_{PQ}
=
\mathcal W_{PQ}\circ \mathcal C_{PQ}\circ \mathcal M_{PQ}\circ \mathcal P_{PQ},
\end{equation}
and an asynchronous update loop that later consensus results can treat as a concrete regulated
model rather than as an abstract rewrite system.
\end{enumerate}
What remains open for \texttt{papers.reality.a.01} is equally important: this note does
\emph{not} prove that the above patch net is already the derived OPH regulator of the compact
paper, does \emph{not} derive the repair law from OPH recovery dynamics, and does \emph{not}
prove the fixed-point hypotheses later needed in the consensus lane. Termination, local
confluence, gauge covariance, and refinement-limit lift remain later tasks.

\subsection{What it feeds now for \texttt{papers.observers.b.03}}\label{sec:microphysics-feeds-b03}

The task \texttt{papers.observers.b.03} asks for the measurement/Born-rule package to be closed.
This note only feeds its input objects; it does not close the package.
\begin{enumerate}[leftmargin=2em]
\item It provides a record-layer candidate: robust pointer observables supported inside patches
and summarized across overlaps.
\item It provides an operational observer criterion
\begin{equation}
O_{\mathrm{ref}}=(P_O,\mathcal A(P_O),\rho_O,R_O,\mathcal U_O),
\end{equation}
which sharpens what fixed-cutoff data the later observer paper would actually need to quantify.
\item It provides simulator diagnostics for record lifetime, repeated-read stability, overlap
comparison, and back-action control, which are the natural phase-1 precursors to a later
measurement theorem package.
\end{enumerate}
What remains open for \texttt{papers.observers.b.03} is the theorem-level measurement work: this
note does \emph{not} prove that the declared record algebra is central or approximately central
in the required sense, does \emph{not} prove definiteness of outcomes, does \emph{not} derive the
Born rule, and does \emph{not} identify collapse with conditioning beyond phase-1 architecture
language. It feeds the measurement package with explicit regulated observables, not with
finished theorems.

\subsection{What it feeds now for \texttt{papers.observers.e.18}}\label{sec:microphysics-feeds-e18}

The task \texttt{papers.observers.e.18} asks for observer continuation / backup to be formalized
as a theorem package. This note again supplies only the fixed-cutoff interface on which that
later package could be built.
\begin{enumerate}[leftmargin=2em]
\item It identifies the checkpoint-like data that could later matter for continuation:
persistent record summaries, overlap-accessible sector data, and local update channels that
preserve those summaries as well as possible under repair.
\item It isolates the collar and patch interfaces through which restoration maps would have to
act, namely the shared observable family on overlaps plus the local record layer inside the
observer patch.
\item It proposes phase-1 diagnostics---record persistence, mismatch decay, and schedule
divergence---that can later become error-accounting inputs for a continuation theorem.
\end{enumerate}
What remains open for \texttt{papers.observers.e.18} is the whole theorem package proper: this
note does \emph{not} define or prove a restoration map, does \emph{not} give accumulated-error
bounds, does \emph{not} state an identity criterion for ``same observer,'' and does \emph{not}
show that repair plus records is sufficient for continuation or backup. It gives the candidate
checkpoint variables and interfaces, not the continuation theorem.

\subsection{Completion-tree status after this note}

The correct completion-tree reading is therefore narrow.
\begin{enumerate}[leftmargin=2em]
\item This note \emph{feeds} \texttt{papers.compact.e.24}, \texttt{papers.reality.a.01},
\texttt{papers.observers.b.03}, and \texttt{papers.observers.e.18} by replacing abstract
placeholders with explicit fixed-cutoff objects.
\item This note does \emph{not} mark any of those leaves complete.
\item The remaining proof burden stays with the later leaf tasks themselves: the weight-law
derivation in \texttt{papers.compact.e.24}, the repair-law theorem chain beginning at
\texttt{papers.reality.a.02}, the measurement/Born-rule closure in
\texttt{papers.observers.b.03}, and the continuation/backup formalization in
\texttt{papers.observers.e.18}.
\end{enumerate}

That is exactly the intended role of an extra phase-1 side paper in the tree: it supplies a
regulated build surface and a reusable object dictionary, while leaving theorem closure and
continuum closure to the papers that already own those tasks.

\section{Conclusion}

The constructive gap addressed by this note is now sharper, but not smaller than it really is.
Imported OPH work already supplies the screen/patch/collar language and the broader overlap
consistency program. The new contribution of this paper is only to make one finite-cutoff
reference architecture explicit enough that patches, overlaps, collar sectors, records,
observer-like subsystems, and synchronization attempts can all be instantiated on the same
regulated screen.

That contribution is substantial at phase 1 because later tasks no longer need to speak only in
generalities. They can refer to explicit link/register Hilbert spaces, explicit overlap
observables, explicit mismatch syndromes, explicit repair primitives, explicit record summaries,
and explicit pilot models. But the note should still be read with discipline. It does \emph{not}
identify the final OPH microphysics, does \emph{not} derive the edge-sector heat-kernel /
Casimir law, does \emph{not} prove the consensus repair theorems, does \emph{not} close the
measurement/Born-rule package, and does \emph{not} prove observer continuation or backup.

So the right final status is strictly phase-1. This note is a regulated architecture surface for
later compact, reality, and observer tasks. If the proposed pilots fail, that falsifies or
deprioritizes this reference architecture. If they succeed, they justify taking the later bridge
tasks seriously on an explicit microscopic model. In either case, the present paper is a build
surface and dependency clarifier, not a theorem-level completion paper.

\end{document}
